%% SoftwareX Original Software Publication (OSP)
%% Based on the official Elsevier template:
%%   https://www.elsevier.com/__data/promis_misc/softwarex-osp-template.tex
%% Local copy: docs/softwarex-osp-template.tex
%% Content describes ctc-discrete-monodromy v0.1.0 (Rust/Python source logic unchanged)

\documentclass[preprint,12pt,a4paper]{elsarticle}

\usepackage{amssymb}
\usepackage{amsmath}
\usepackage{hyperref}
\usepackage{float}
\restylefloat{table}
\usepackage{booktabs}
\usepackage{graphicx}
\usepackage{listings}
\setlength{\parindent}{0pt}

\lstset{
  basicstyle=\ttfamily\footnotesize,
  breaklines=true,
  frame=single,
  columns=fullflexible,
  keepspaces=true,
  showstringspaces=false
}

\newcommand{\manuscriptdate}{23 July 2026}
\newcommand{\Phir}{\Phi(r)}
\newcommand{\br}{b(r)}
\newcommand{\omegar}{\omega(r)}
\newcommand{\MN}{\mathcal{M}_{N}}

\journal{SoftwareX}

\begin{document}
\renewcommand{\labelenumii}{\arabic{enumi}.\arabic{enumii}}

\begin{frontmatter}

\title{ctc-discrete-monodromy: a reproducible Rust/SymPy framework for discrete operator diagnostics on axisymmetric ergoregion backgrounds}

\author[umsn]{Mauricio Miguel Paniagua Ramirez\corref{cor1}}
\ead{1912574a@umich.mx}
\cortext[cor1]{Corresponding author}
\address[umsn]{Facultad de Ingenier\'ia El\'ectrica,
Universidad Michoacana de San Nicol\'as de Hidalgo (UMSNH),
Morelia, Michoac\'an, M\'exico}

\begin{abstract}
\texttt{ctc-discrete-monodromy} is an open-source Rust/SymPy research-software stack for discrete operator diagnostics on explicit axisymmetric metric profiles: circulant-kernel assembly, spectral-radius and Lipschitz regression oracles, homogeneous and inhomogeneous Picard solvers ($\Phi=K\Phi+S$), modular kernel comparison, and RK4 geodesic checks.
Although motivated by frame-dragging Morris--Thorne backgrounds, the numerical core is domain-transferable to any periodic sampling loop that needs reproducible row-normalized dense maps and fixed-point residual gates.
Design choices such as mass-gap row renormalization are implemented as software test oracles ($\rho(K)<1$, residuals $<10^{-8}$), not continuum field theorems.
All $17/17$ Rust unit tests pass; the stack is MIT-licensed and runs on Windows, Linux, and macOS.
\end{abstract}

\begin{keyword}
scientific software \sep Rust \sep SymPy \sep discrete operators \sep spectral diagnostics \sep reproducible research
\end{keyword}

\end{frontmatter}

%%==============================================================================
\section*{Required Metadata}
\label{sec:required-metadata}

\section*{Current code version}
\label{sec:code-version}

Ancillary data table required for subversion of the codebase.

\begin{table}[H]
\small
\begin{tabular}{|l|p{5.2cm}|p{6.8cm}|}
\hline
\textbf{Nr.} & \textbf{Code metadata description} & \textbf{Please fill in this column} \\
\hline
C1 & Current code version & 0.1.0 \\
\hline
C2 & Permanent GitHub link to code/repository used for this code version & \url{https://github.com/mauramyrez/ctc-discrete-monodromy} \\
\hline
C3 & Legal Code License & MIT License (\texttt{LICENSE} / \texttt{LICENSE.txt} in repository root) \\
\hline
C4 & Code versioning system used & git (annotated tag \texttt{v0.1.0}) \\
\hline
C5 & Software code languages, tools, and services used & Rust 2021 edition (\texttt{nalgebra}, \texttt{serde}, \texttt{anyhow}, \texttt{thiserror}); optional Python~3 (SymPy, NumPy, SciPy, Matplotlib) for symbolic export and figure regeneration; Cargo; pip \\
\hline
C6 & Compilation requirements, operating environments \& dependencies & \texttt{rustup} (stable) + Cargo; optional \texttt{high-prec} feature (GMP/MPFR via \texttt{rug}). OS: Windows~10+, Linux, macOS. Optional Python~$\ge$3.10 with \texttt{requirements.txt} pins (not required for \texttt{cargo test}). \\
\hline
C7 & If available Link to developer documentation/manual & \url{https://github.com/mauramyrez/ctc-discrete-monodromy\#readme} \\
\hline
C8 & Support email for questions & 1912574a@umich.mx \\
\hline
\end{tabular}
\caption{Code metadata (mandatory for SoftwareX OSP submissions)}
\label{tab:code-meta}
\end{table}

%%==============================================================================
\section{Motivation and significance}
\label{sec:motivation}

Research software for continuum tensor calculus and spacetime visualization is comparatively mature~\cite{Meurer2017,Bapat2019,SageManifolds,BaumgarteShapiro2010}, while reproducible hygiene for dense circulant operators---row-sum normalization semantics, spectral-radius regression, Picard residual tracking, and dimension-safe matrix--vector application---is still often left to one-off notebooks built on generic array stacks~\cite{Harris2020,Virtanen2020,Trefethen2000}.
That gap produces fragile floating-point pipelines that are difficult to cite, fork, and regress across operating systems.

\texttt{ctc-discrete-monodromy} addresses this gap as a typed Rust engine with optional SymPy documentation exports.
The historical package name retains the prefix \texttt{ctc-} for repository continuity; the software does \emph{not} claim continuum closed-timelike-curve physics.
The reference application samples an explicit Morris--Thorne profile class with azimuthal frame dragging~\cite{Morris1988,Visser1995,HawkingEllis1973}, but the reusable product is the discrete-operator harness: any user who needs a maintainable benchmark for row-normalized maps on periodic nodes can adopt the same APIs, tests, and CLI reports.

An open-source Rust implementation is necessary because:
\begin{itemize}
  \item typed ownership and explicit dimension asserts eliminate silent truncation bugs common in notebook-centric operator code;
  \item release-mode performance with link-time optimization supports dense $N\times N$ kernel builds and Picard loops without sacrificing reproducibility;
  \item a Cargo test harness provides deterministic regression gates ($17/17$ unit tests) suitable for continuous integration and cross-domain reuse studies.
\end{itemize}

\paragraph{Experimental setting (how the user uses the software).}
A typical session is:
(1)~build the Rust crate and run \texttt{ctc-engine} for interactive spectral/Picard reports;
(2)~execute \texttt{cargo test} as the acceptance gate before modifying kernels or profiles;
(3)~optionally install the pinned Python stack to regenerate SymPy tensor assets or illustration figures.
All certified numerical claims refer to the discrete operator $\MN$ constructed in the crate; continuous equations only motivate kernel shapes.

\paragraph{Related work.}
Symbolic and numerical GR ecosystems emphasize continuum tensors and spacetime analysis~\cite{Meurer2017,Bapat2019,SageManifolds,BaumgarteShapiro2010}.
Spectral and array libraries provide the generic linear-algebra substrate~\cite{Trefethen2000,Harris2020,Virtanen2020} but do not ship a regression-tested, dimension-safe circulant monodromy oracle with inhomogeneous Picard workloads and modular kernel comparison.
This software fills that narrow but reusable niche and archives the implementation under an MIT license~\cite{PaniaguaRamirez2026CTC}.

%%==============================================================================
\section{Software description}
\label{sec:software}

\subsection{Software architecture}
\label{sec:architecture}

The GitHub repository root is the Cargo package root (Fig.~\ref{fig:pipeline}):

\begin{lstlisting}[language={}]
ctc-discrete-monodromy/          # repository / crate root
  Cargo.toml, LICENSE, LICENSE.txt, README.md, requirements.txt
  src/                           # Rust crate: ctc_research_framework
    lib.rs                       # algebra, dynamics, energy
    algebra/                     # metric profiles, invariants, Christoffel FD
    dynamics/
      monodromy.rs               # DiscreteMonodromyOperator, Picard, kernels
      geodesic.rs                # RK4 equatorial geodesic integrator
      bvp.rs                     # smoke tests (homogeneous + inhomogeneous)
    energy/                      # API scaffold for quadratic-form checks
    main.rs                      # ctc-engine CLI driver
    tests/monodromy_test.rs      # regression suite entry
  scripts/                       # optional SymPy export + Matplotlib figures
  docs/                          # SoftwareX OSP manuscript + figures
\end{lstlisting}

\begin{itemize}
  \item \textbf{\texttt{lib.rs}} exposes \texttt{algebra}, \texttt{dynamics}, and \texttt{energy}, with a \texttt{[cfg(test)]} suite entry point.
  \item \textbf{\texttt{monodromy.rs}} is the core deliverable: \texttt{PeriodicField}, \texttt{DiscreteMonodromyOperator}, modular \texttt{KernelKind} assembly, row-sum renormalization, power-iteration spectral diagnostics, grid-oscillation Lipschitz bounds, and Picard iteration for both homogeneous ($S=0$) and inhomogeneous ($\Phi=K\Phi+S$) maps.
  \item \textbf{SymPy scripts} (\texttt{derive\_tensors.py}, \texttt{verify\_energy\_conditions.py}) generate static equatorial tensor assets; they are not invoked from the Rust hot path.
  \item \textbf{\texttt{ctc-engine}} (\texttt{main.rs}) prints geometry, $\rho(K)$, oscillation Lipschitz, homogeneous/inhomogeneous Picard residuals, modular-kernel comparison, and energy-scaffold smoke results.
\end{itemize}

\begin{figure}[t]
\centering
\fbox{\parbox{0.92\linewidth}{\centering\footnotesize\sffamily
\textbf{SymPy} (\texttt{scripts/})\;$\longrightarrow$\;
\textbf{Rust crate} (\texttt{lib.rs} / \texttt{monodromy.rs})\;$\longrightarrow$\;
\textbf{cargo test} (17/17)\;$\longrightarrow$\;
\textbf{docs/} (SoftwareX \LaTeX)
}}
\caption{End-to-end software pipeline: optional symbolic export, typed discrete-operator engine, automated regression, and SoftwareX manuscript assets.}
\label{fig:pipeline}
\end{figure}

\paragraph{Reference geometry configuration.}
Profiles used throughout the crate are closed-form,
\begin{equation}
\label{eq:profiles}
\br = r_{0}\Bigl(\frac{r_{0}}{r}\Bigr)^{\gamma},\quad
\Phir = -\frac{\alpha r_{0}}{r},\quad
\omegar = \omega_{0}\exp\!\Bigl[-\beta\frac{r-r_{0}}{r_{0}}\Bigr],
\end{equation}
with defaults $(r_{0},\gamma,\alpha,\omega_{0},\beta)=(1.0,\,0.5,\,0.1,\,1.2,\,2.0)$
via \texttt{ExplicitProfileParams::prd\_reference()}.
The monodromy loop defaults to $r_{\mathrm{loop}}=1.5\,r_{0}$ (just outside the reported outer ergosurface $r_{\mathrm{ERGO}}\approx 1.238\,r_{0}$), providing a strongly dragged near-ergoregion exterior without throat-adjacent coordinate singularities.

\subsection{Software functionalities}
\label{sec:functionalities}

\paragraph{Discrete monodromy operator (core API).}
On $N$ uniform azimuthal nodes the crate builds a sampled proxy matrix $\widetilde{K}$ (modular \texttt{KernelKind}) and applies an \emph{explicit} row-sum renormalization
\begin{equation}
\label{eq:K-renorm}
K_{ij}
=
\widetilde{K}_{ij}\,
\frac{\exp(-m_{\ast}\Delta\tau)}{\sum_{k}\widetilde{K}_{ik}},
\end{equation}
producing the circulant (constant-row-sum) operator $\MN:\mathbb{R}^{N}\to\mathbb{R}^{N}$.
Row renormalization is a software design pattern and numerical test oracle: it forces every row sum to $\alpha=\exp(-m_{\ast}\Delta\tau)<1$ whenever $m_{\ast}>0$, hence $\rho(K)=\alpha<1$ by construction for nonnegative constant-row-sum matrices.
The homogeneous map $\Phi\mapsto K\Phi$ therefore converges to the trivial fixed point $\Phi^{*}=0$, which the regression suite uses as a floating-point stress test---not as a continuum Green-function solution.

\paragraph{Inhomogeneous response and modular kernel diagnostics.}
\label{sec:inhomogeneous}
Beyond the homogeneous oracle, the solver implements the discrete inhomogeneous fixed-point map
\begin{equation}
\label{eq:inhomo}
\Phi = K\Phi + S,
\end{equation}
equivalently $\Phi=(I-K)^{-1}S$ when $\rho(K)<1$.
Picard iteration advances $\Phi_{n+1}=K\Phi_n+S$ and reports the residual $\|\Phi-K\Phi-S\|_{2}$.
A default sinusoidal forcing vector $S$ (unit-scaled mode on the azimuthal grid) produces a nontrivial response $\|\Phi^{*}\|_{2}>0$, elevating the computational workload above a pure null-solution regression.
Separately, \texttt{KernelKind::\{YukawaProxy, ExponentialDecay\}} exposes comparative modular kernel assembly under identical row-renormalization semantics, so spectral-radius and Lipschitz oracles can be contrasted across kernel families without changing the metric profile or node count.
Both paths are discrete operator diagnostics on explicit metric profiles; they do not claim continuum PDE well-posedness.

\paragraph{Spectral and Lipschitz diagnostics.}
\texttt{DiscreteMonodromyOperator::spectral\_radius} (power iteration) and \texttt{lipschitz\_bound}
\[
L\le\rho(K)\,(1+\pi\sqrt{2}/N)
\]
are exposed as software thresholds.
\texttt{PeriodicField} provides $L^{2}$ and index-space oscillation norms; \texttt{apply\_kernel} asserts \texttt{field.len()==n\_nodes} to forbid silent truncation.

\paragraph{Geodesic integrator.}
\texttt{geodesic.rs} supplies classical RK4 integration ($\Delta\tau=0.01$, $200$ steps) of co-rotating, unit-normalized equatorial initial data with Christoffel symbols from central finite differences ($\varepsilon=10^{-5}$).
These orbits are integrator regression checks on the near-ergoregion exterior.

\paragraph{Optional symbolic export and illustration workflow.}
Python scripts optionally export equatorial Einstein-tensor components and NEC/WEC contractions for the profile class~\eqref{eq:profiles}, and can regenerate vector illustration PDFs under \texttt{docs/figures/}.
These helpers support manuscript documentation; they are not required to build or regress the Rust engine.
Rust \texttt{energy} modules exercise only scaffold quadratic forms for API sanity; manuscript-grade contractions remain SymPy assets when the optional Python layer is used.

\subsection{Sample code snippets analysis}
\label{sec:snippets}

The default factory \texttt{default\_monodromy\_operator()} wires the reference profile at $r_{\mathrm{loop}}=1.5\,r_{0}$ with $m^{2}=1$, $\nu=1$, and $N=128$.
\texttt{find\_fixed\_point} (homogeneous) and \texttt{find\_fixed\_point\_inhomogeneous} rebuild $K$ once, then Picard-iterate with residual tolerance $10^{-12}$ (solver) under the contraction gate \texttt{is\_contraction()}.
Dimension mismatches panic early, which is intentional software hygiene for regression oracles.

%%==============================================================================
\section{Illustrative examples}
\label{sec:examples}

\subsection{Build, run, and reproduce the test suite}

\begin{lstlisting}[language={}]
# Rust discrete-operator engine (acceptance gate)
# Optional on Windows long paths: set CARGO_TARGET_DIR=C:\ctc-target
cargo build --release
cargo run --release --bin ctc-engine
cargo test   # expect 17/17

# Optional Python symbolic / figure layer (documentation only)
python -m venv .venv
# Windows: .venv\Scripts\activate
pip install -r requirements.txt
python scripts/derive_tensors.py
python scripts/plot_spacetime_diagrams.py
\end{lstlisting}

Expected outcome of \texttt{cargo test}: \textbf{17/17} unit tests passed, covering
(i)~metric / flare-out / ergosurface diagnostics;
(ii)~energy-scaffold API checks;
(iii)~monodromy bounds $\rho(K)<1$ and $L<1$;
(iv)~homogeneous Picard residuals $<10^{-8}$ in $<500$ iterations for $N\in\{32,64,128\}$;
(v)~inhomogeneous response $\Phi=K\Phi+S$ with nontrivial $\|\Phi^{*}\|_{2}$ and modular-kernel spectral comparison;
(vi)~RK4 co-rotating geodesic integration.

\subsection{Constructing and diagnosing $\MN$ in Rust}

\begin{lstlisting}[language={}]
use ctc_research_framework::dynamics::monodromy::{
    default_forcing_vector, default_monodromy_operator,
    find_fixed_point, find_fixed_point_inhomogeneous, PeriodicField,
};

let mut m = default_monodromy_operator(); // N=128, m^2=1, nu=1
m.n_nodes = 64;
let init = PeriodicField::new(
    (0..m.n_nodes).map(|i| 1.0 + 0.1 * (i as f64)).collect(),
);
assert!(m.is_contraction());           // rho(K)<1 and L<1
let fp0 = find_fixed_point(&init, &m)?; // homogeneous -> Phi* ~ 0
assert!(fp0.field.l2_norm() < 1e-6);

let source = default_forcing_vector(m.n_nodes);
let zero = PeriodicField::new(vec![0.0; m.n_nodes]);
let fp = find_fixed_point_inhomogeneous(&zero, &m, &source)?;
assert!(fp.residual < 1e-8);
assert!(fp.field.l2_norm() > 1e-3);    // nontrivial response
\end{lstlisting}

Default configuration parameters are summarized in Table~\ref{tab:params}.
Measured discrete diagnostics for the reference profile appear in Table~\ref{tab:monodromy}.

\begin{table}[H]
\centering
\caption{Default configuration parameters for the monodromy regression baseline.}
\label{tab:params}
\begin{tabular}{@{}ll@{}}
\toprule
Parameter & Default value \\
\midrule
Profile & \texttt{ExplicitProfileParams::prd\_reference()} \\
$r_{\mathrm{loop}}$ & $1.5\,r_{0}$ \\
$m^{2}$, $\nu$, $N$ & $1$, $1$, $128$ \\
Solver tolerance & $10^{-12}$ \\
Test residual threshold & $10^{-8}$--$10^{-10}$ \\
Picard iteration cap & $500$ (tests) / $10^{4}$ (solver) \\
RK4 step / steps & $\Delta\tau=0.01$ / $200$ \\
\bottomrule
\end{tabular}
\end{table}

\begin{table}[H]
\centering
\caption{Discrete monodromy diagnostics at $r_{\mathrm{loop}}=1.5\,r_{0}$ ($m^{2}=1$, $\nu=1$).
$\rho(K)$ via power iteration; $L$ is the grid-oscillation Lipschitz diagnostic; residuals are $\|\Phi_{k}-K\Phi_{k}\|_{2}$ after homogeneous Picard iteration (baseline). Inhomogeneous runs with default forcing $S$ yield nontrivial $\|\Phi^{*}\|_{2}$ at comparable residual tolerance.}
\label{tab:monodromy}
\begin{tabular}{@{}ccccc@{}}
\toprule
$N$ & $\rho(K)$ & $L$ & residual & Picard iterations \\
\midrule
$32$ & $0.702$ & $0.799$ & $7.8\times 10^{-13}$ & $79$ \\
$64$ & $0.838$ & $0.896$ & $9.8\times 10^{-13}$ & $155$ \\
$128$ & $0.915$ & $0.947$ & $9.9\times 10^{-13}$ & $308$ \\
\bottomrule
\end{tabular}
\end{table}

\subsection{CLI output shape}

Running \texttt{cargo run --release --bin ctc-engine} prints a geometry line ($r_{\mathrm{ERGO}}$, $g_{tt}(r_{\mathrm{loop}})$), a homogeneous Picard table with columns
\texttt{N}, \texttt{rho(K)}, \texttt{L\_osc}, \texttt{residual}, \texttt{iters}, \texttt{||Phi*||\_2} (matching Table~\ref{tab:monodromy} to working precision), an inhomogeneous $\Phi=K\Phi+S$ residual line, and a modular-kernel $\rho(K)$ comparison.
Figure scripts under \texttt{scripts/plot\_spacetime\_diagrams.py} optionally regenerate ergosurface, light-cone tipping, and geodesic-loop diagrams into \texttt{docs/figures/} for manuscript illustration.

%%==============================================================================
\section{Impact}
\label{sec:impact}

The primary impact of \texttt{ctc-discrete-monodromy} is reusable research-software infrastructure for discrete linear maps, not a claim of new physical law.
SoftwareX-relevant reuse pathways include:

\begin{itemize}
  \item \textbf{Cross-domain numerical analysis.}
  Any periodic sampling problem that needs a row-normalized dense kernel, spectral-radius / Lipschitz gates, and Picard residual tracking can reuse \texttt{DiscreteMonodromyOperator}, \texttt{KernelKind}, and $\Phi=K\Phi+S$ without adopting the Morris--Thorne profile layer.
  \item \textbf{Computational-relativity workflows.}
  Relative to continuum GR packages~\cite{Bapat2019,SageManifolds,BaumgarteShapiro2010}, the crate supplies a missing regression harness for discrete circulant diagnostics on explicit axisymmetric backgrounds~\cite{PaniaguaRamirez2026CTC}.
  \item \textbf{Daily practice for research software engineers.}
  A single command path (\texttt{cargo test} / \texttt{ctc-engine}) turns profile construction into spectral/Picard reports with dimension-safe asserts, reducing notebook fragility.
  \item \textbf{Benchmarking and CI.}
  Circulant row-normalization and oscillation Lipschitz bounds act as portable oracles: kernel variants that break $\rho(K)<1$ or residual thresholds fail loudly under the $17/17$ suite.
  \item \textbf{Educational reuse.}
  The optional SymPy/Matplotlib layer regenerates illustration assets for teaching frame-dragging geometry while keeping the Rust acceptance gate independent of Python.
\end{itemize}

As an initial public release (v0.1.0), download and citation metrics are not yet available; the near-term impact target is archival citability and adoption as a small, forkable discrete-operator benchmark within scientific computing and numerical-relativity communities.

%%==============================================================================
\section{Conclusions}
\label{sec:conclusions}

We have delivered an open-source Rust/SymPy framework whose reusable core is a regression-tested discrete operator engine---row-normalized circulant kernels, spectral/Lipschitz diagnostics, homogeneous and inhomogeneous ($\Phi=K\Phi+S$) Picard solvers, and modular kernel comparison---demonstrated on explicit frame-dragging Morris--Thorne profiles with optional symbolic export.
The contribution is a high-performance, reproducible software artifact aligned with SoftwareX's emphasis on reusable research software, with clear separation between implemented numerics and analytical outlook.
Source code, license (MIT), and documentation are permanently available at
\url{https://github.com/mauramyrez/ctc-discrete-monodromy}~\cite{PaniaguaRamirez2026CTC}
(manuscript date: \manuscriptdate).

\section*{Acknowledgements}
\label{sec:ack}
Symbolic and numerical artifacts accompany this SoftwareX Original Software Publication in the open-source repository \texttt{ctc-discrete-monodromy}.

\bibliographystyle{elsarticle-num}
\bibliography{references}

\section*{Current executable software version}
\label{sec:soft-version}

\begin{table}[H]
\small
\begin{tabular}{|l|p{5.2cm}|p{6.8cm}|}
\hline
\textbf{Nr.} & \textbf{(Executable) software metadata description} & \textbf{Please fill in this column} \\
\hline
S1 & Current software version & 0.1.0 (\texttt{ctc-research-framework} / \texttt{ctc-engine}) \\
\hline
S2 & Permanent link to executables of this version & \url{https://github.com/mauramyrez/ctc-discrete-monodromy} \\
\hline
S3 & Legal Software License & MIT License \\
\hline
S4 & Computing platforms / Operating Systems & Windows, Linux, macOS \\
\hline
S5 & Installation requirements \& dependencies & See Table~\ref{tab:code-meta} (C5--C6); \texttt{cargo build --release}; \texttt{pip install -r requirements.txt} \\
\hline
S6 & If available, link to user manual---if formally published include a reference to the publication in the reference list & Repository README; this SoftwareX article \\
\hline
S7 & Support email for questions & 1912574a@umich.mx \\
\hline
\end{tabular}
\caption{Software metadata (optional)}
\label{tab:soft-meta}
\end{table}

\end{document}
